\SourceLeaf{20}

Reversing the preceding derivation shows that from \textsc{Gauss}'s statements, in the
case $q=3n+1$, one can first infer the proportion
\[
\left[\dfrac{\varkappa_1}{\pi_1}\right]:\left[\dfrac{\pi_1}{\varkappa_1}\right]=\left[\dfrac{\varkappa_1}{\pi_2}\right]:\left[\dfrac{\pi_2}{\varkappa_1}\right]
\]
Applying once more the law expressed in this proportion, one finds
\[
\dfrac{\left[\dfrac{\varkappa_1}{\pi_1}\right]}{\left[\dfrac{\pi_1}{\varkappa_1}\right]}=\dfrac{\left[\dfrac{\varkappa_2}{\pi_1}\right]}{\left[\dfrac{\pi_1}{\varkappa_2}\right]}=\dfrac{\left[\dfrac{\varkappa_1}{\pi_2}\right]^2}{\left[\dfrac{\pi_2}{\varkappa_1}\right]^2}
\quad\text{and hence}\quad
\dfrac{\left[\dfrac{\varkappa_1}{\pi_2}\right]}{\left[\dfrac{\pi_2}{\varkappa_1}\right]}=\dfrac{\left[\dfrac{\varkappa_1}{\pi_2}\right]^3}{\left[\dfrac{\pi_2}{\varkappa_1}\right]^2}.
\]
Since $\left[\dfrac{\varkappa_1}{\pi_2}\right]=\left[\dfrac{\pi_2}{\varkappa_1}\right]$ follows from this, the general reciprocity law in its simplest form is, for $q=3n+1$ as well, implicitly contained in \textsc{Gauss}'s statements.

The notes in which \textsc{Gauss} proved his cubic reciprocity theorems are collected under [4]. Since these notes come from several undated loose slips, the date of the proofs in question cannot be determined with certainty. It is
likely, however, that the derivations were carried out soon after he obtained the »Observatio« and the »Criterien«,
for their line of reasoning is very closely connected with the first statements under [2]
and, moreover, with Art. 358 of the Disquisitiones arithmeticae. On the other hand, \textsc{Gauss}, by his own
account\HistoricalNote{\textsuperscript{*)}}{*)}{See Art. 33 of the treatise »Summatio quarumdam serierum singularium«, presented to the Göttingen Society on 24 August 1808.}, was in any case in possession, before the middle of 1808, of the sixth proof of the quadratic
reciprocity law. That proof is the one most closely related to the proofs under consideration here.

\textsc{Gauss}'s proof of the cubic reciprocity law is based on cyclotomy. The
individual formulas under [4] can either be taken directly from Art. 358 of the Disquisitiones arithmeticae or at least be derived with very little effort from the developments given there.
The congruence
\[
P^{pp}\equiv r^{pp}+Rr^{ppg}+\ldots\equiv P(R^{\mathrm{ind.}\,p})
\]
is modulo $p$ and arises from the fact that all polynomial coefficients of the power $pp$ other than 1 are divisible by $p$. The distinction between the cases
\[
\left\{\dfrac{M-N\sqrt{-27}}{M+N\sqrt{-27}}\right\}^{\frac13(pp-1)}\equiv 1,\quad -\tfrac12+\tfrac12\sqrt{-3},\quad -\tfrac12-\tfrac12\sqrt{-3}\qquad\;(\mathrm{mod.}\ p)
\]
corresponds exactly, for both $p=3m-1$ and $p=3m+1$, to the distinction made above under I, II and III,
namely $B\equiv 0\;(\mathrm{mod.}\ q)$, etc. The correspondence between those three congruences and the following
\[
p^{\frac13(n-1)}\equiv 1,\quad -\tfrac12-\tfrac12\dfrac{M}{3N},\quad -\tfrac12+\tfrac12\dfrac{M}{3N}\qquad\;(\mathrm{mod.}\ n),
\]
as developed in [4], therefore already contains a complete proof of \textsc{Gauss}'s statements
and thus implicitly of the general cubic reciprocity law.

The further discussion under [4] concerns the criteria for 2 and 3 to be cubic residues.
As for 3 as a cubic residue, it suffices to note that (using the notation in [4])
$p^3\equiv p, p'$ or $p''\;(\mathrm{mod.}\ 3)$ holds according to whether ind.\,$3\equiv 0, 1$ or $2\;(\mathrm{mod.}\ 3)$ holds. The condition
$N\equiv 0\;(\mathrm{mod.}\ 3)$ for 3 to be a cubic residue modulo $n$ then follows directly from the concluding statements of [4].

